\documentclass[11pt]{article}
\usepackage[utf8]{inputenc}
\usepackage{amsmath,amssymb,amsthm}
\usepackage[margin=1in]{geometry}
\usepackage{hyperref}
\usepackage{xcolor}
\usepackage{enumitem}
\usepackage{newunicodechar}
\newunicodechar{ℝ}{\ensuremath{\mathbb{R}}}
\newunicodechar{⟨}{\ensuremath{\langle}}
\newunicodechar{⟩}{\ensuremath{\rangle}}
\newunicodechar{κ}{\ensuremath{\kappa}}
\newunicodechar{φ}{\ensuremath{\varphi}}
\newunicodechar{μ}{\ensuremath{\mu}}
\newunicodechar{λ}{\ensuremath{\lambda}}
\newunicodechar{α}{\ensuremath{\alpha}}
\newunicodechar{γ}{\ensuremath{\gamma}}
\newunicodechar{ε}{\ensuremath{\varepsilon}}

\newcommand{\userbox}[1]{\par\medskip\begin{quote}\small\textbf{\textcolor{blue!60}{DM:}} #1\end{quote}\medskip}
\newcommand{\claudebox}[1]{\par\medskip\begin{quote}\small\textbf{\textcolor{orange!60!black}{Claude:}} #1\end{quote}\medskip}
\newcommand{\gptbox}[1]{\par\medskip\begin{quote}\small\textbf{\textcolor{green!50!black}{GPT-5.5 Pro:}} #1\end{quote}\medskip}

\title{Tide retrospective:\\concrete \texttt{GibbsRegularity} for the harmonic Gibbs}
\author{Daniel Murfet}
\date{6 May 2026}

\begin{document}
\maketitle

\begin{abstract}
This retrospective records Tide step 10, the second tide of the day,
which provides a concrete \texttt{GibbsRegularity} instance for the
1D harmonic potential $L(x) = (\lambda/2)x^2$ with the linear
perturbation $A(x) = x$ on $\mathbb{R}$ ($\lambda, t > 0$). The
instance unlocks threepoint's \emph{derivative} theorems --- the FDT,
cross-susceptibility, and flow equation --- for the harmonic case,
transforming them from parameterised lemmas (conditional on a
hypothesis bundle) into theorems about the harmonic Gibbs measure
proper. The formalisation took 265 lines of Lean in a new file
\texttt{Laplace/OneD/HarmonicGibbsRegularity.lean}, with zero
\texttt{sorry} / \texttt{axiom} / \texttt{native\_decide}, and landed
in a single session. The single decisive contribution from the
GPT-5.5~Pro deliberation was a strategic pivot away from the
expected dominated-convergence approach
(\texttt{hasDerivAt\_integral\_of\_dominated\_loc\_*}) toward a
\emph{closed-form} proof of the partition function for all $h$,
which sidesteps the differentiation-under-the-integral machinery
entirely and produces a stronger primitive theorem as a free side
effect.
\end{abstract}

\section{Setting}

The kappa3-harmonic tide (Tide 5) and the anharmonic-$\kappa_3$ tide
(Tide 9) both deliberately avoided \texttt{GibbsRegularity}. The
reason was given in Tide 5's GPT consult, quoted verbatim in that
tide's retrospective: \emph{``\texttt{kappa3} is defined using
\texttt{gibbsExp~...~0~...}, i.e.\ at $h = 0$. To show
$\kappa_3(x,x,x) = 0$ you only need $\langle x \rangle = 0$,
$\langle x^3 \rangle = 0$. You do not need the FDT differentiation
machinery.''} Both tides therefore proved their cumulant statements
without instantiating the regularity bundle.

The cost of that decision was that threepoint's \emph{derivative}
theorems --- the FDT
($\partial_h \langle \varphi \rangle_h\!\restriction_{h=0}
= -t \cdot \mathrm{Cov}(\varphi, A)$),
the cross-susceptibility identity
($\partial_h \mathrm{Cov}_h(\varphi, B)\!\restriction_{h=0}
= -t \cdot \kappa_3(\varphi, A, B)$),
and the flow equation --- remained \emph{parameterised} on the
hypothesis bundle. They are theorems of the form ``if you give me a
\texttt{GibbsRegularity} instance for $(L, A, t)$, here's the FDT''
rather than ``here's the FDT for the harmonic Gibbs.''

The 6~May candidates survey listed providing the harmonic instance
as \textbf{C3}: ``the natural follow-up'' to Tide 5, with an estimated
size of 150--250 lines. The expected approach was to apply
Mathlib's \texttt{hasDerivAt\_integral\_of\_dominated\_loc\_of\_lip}
(or its derivative-bound variant) with a square-completion
dominator on a neighbourhood of $h = 0$.

\section{The thing formalised}

The headline statement is:

\begin{quote}\itshape
For $\lambda > 0$ and $t > 0$, the triple
$(\mu, L, A) = (\mathrm{volume},\, (\lambda/2)x^2,\, x)$ on $\mathbb{R}$
satisfies the regularity bundle:
\[
  \mathtt{Threepoint.GibbsRegularity}\;
    (\mathrm{volume})\; (\lambda/2 \cdot \square^2)\; (\mathrm{id})\; t.
\]
\end{quote}

The Lean signature:\footnote{File:
\texttt{Laplace/OneD/HarmonicGibbsRegularity.lean} on
\texttt{tide/harmonic-gibbsregularity} at commit \texttt{3e95014}.}

\begin{verbatim}
theorem Threepoint.harmonic_id_gibbsRegularity
    {lam t : R} (hlam : 0 < lam) (ht : 0 < t) :
    Threepoint.GibbsRegularity (volume : Measure R)
      (fun x : R => lam / 2 * x ^ 2)
      (fun x : R => x) t
\end{verbatim}

(Writing \texttt{R} for the type of real numbers, which the Lean
source spells \texttt{$\mathbb{R}$}.)

The instance discharges three fields: strict positivity of the
unperturbed partition, the $h = 0$ identity (a definitional
simp-reduction), and differentiability of the partition under the
integral at $h = 0$ with the explicit pointwise-derivative integral
as the derivative value.

The headline theorem is named in the \texttt{Threepoint} namespace
(via the \texttt{\_root\_} prefix\footnote{Lean syntax:
\texttt{\_root\_.Threepoint.harmonic\_id\_gibbsRegularity}.}) even
though it lives in the laplace repo. This is consistent with the
cross-repo architecture established by Tide~9 and recommended by GPT:
laplace depends on threepoint, and the harmonic instance for
threepoint's structure is naturally proven where laplace's
asymptotic machinery is available, while remaining
namespace-discoverable as a threepoint-side fact.

\section{Proof strategy}

The proof structure is shorter than the dominated-convergence
approach would have been. The pivot to the closed-form route was
the decisive Step~2 contribution.

\paragraph{The closed form.} For any $h \in \mathbb{R}$ and
$\lambda, t > 0$,
\[
  Z(h) := \int_{\mathbb{R}} e^{-t\,((\lambda/2)x^2 + h\,x)}\,dx
  \;=\; e^{t h^2 / (2\lambda)} \sqrt{2\pi / (\lambda t)}.
\]
The proof is a four-step calculation:
\begin{enumerate}
\item \emph{Square completion}, as a pointwise identity:
\[
  (\lambda/2) x^2 + h\,x
  = (\lambda/2) (x + h/\lambda)^2 - h^2 / (2\lambda).
\]
\item \emph{Boltzmann factor decomposition}:
\[
  e^{-t((\lambda/2)x^2 + h\,x)}
  = e^{t h^2 / (2\lambda)} \cdot e^{-(t (\lambda/2)) (x + h/\lambda)^2}.
\]
\item \emph{Translation invariance} of Lebesgue measure
(\texttt{integral\_add\_right\_eq\_self}, generated by
\texttt{to\_additive} from \texttt{integral\_mul\_right\_eq\_self}):
\[
  \int e^{-(t (\lambda/2)) (x + h/\lambda)^2}\,dx
  = \int e^{-(t (\lambda/2)) y^2}\,dy.
\]
\item \emph{Mathlib's} \texttt{integral\_gaussian} at $b = t\lambda/2$:
\[
  \int e^{-b y^2}\,dy = \sqrt{\pi/b}
  \;=\; \sqrt{2\pi/(\lambda t)}.
\]
\end{enumerate}

\paragraph{The three regularity fields, from the closed form.}
\begin{itemize}
\item \emph{Partition positivity.} $Z(0) = \sqrt{2\pi/(\lambda t)}$, and
$\sqrt{\cdot}$ of a positive real is positive (\texttt{Real.sqrt\_pos}).

\item \emph{The $h = 0$ identity.}
$\int e^{-t((\lambda/2)x^2 + 0\cdot x)}\,dx
= \int e^{-t(\lambda/2)x^2}\,dx$ by integrand congruence (the
\texttt{0~*~x} simplifies to \texttt{0}, and \texttt{(...)~+~0~=~(...)}).

\item \emph{Differentiability at $h = 0$.} The partition function in
$h$ equals the closed-form expression
$h \mapsto e^{t h^2/(2\lambda)} \cdot \sqrt{2\pi/(\lambda t)}$. The
derivative of this expression at $h = 0$ is zero, since
$\frac{d}{dh}[t h^2/(2\lambda)] = t h / \lambda$ vanishes at $h = 0$,
and $\exp(0) \cdot 0 = 0$. The target derivative value
$\int (-t \cdot x) e^{-(t \lambda/2) x^2}\,dx$ is also zero by
parity. So both sides equal zero, and \texttt{HasDerivAt} closes by
the closed-form differentiation.
\end{itemize}

The mathematical content is short. Most of the 265 lines are Lean
plumbing: rewriting integrand lambdas, congruence lemmas, the
\texttt{HasDerivAt.comp} chain for the closed-form derivative.

\section{Roadblocks and resolutions}

\paragraph{The strategic pivot.} The pre-deliberation plan was to
apply Mathlib's dominated-convergence theorem for differentiation
under the integral, with a square-completion
dominator.\footnote{The two relevant Mathlib lemmas live in
\texttt{ParametricIntegral.lean}: one takes a Lipschitz bound, the
other takes a derivative bound. Either suffices for the harmonic
case.} GPT
named this \emph{Candidate A} and replaced it with \emph{Candidate
A$'$}: the closed-form route. The verbatim recommendation:

\gptbox{Short version: I would \textbf{not} lead with dominated
differentiation for this instance. For the concrete harmonic+linear
case, the cleanest proof is usually: (1) prove the closed form
$Z(h) = e^{t h^2/(2\lambda)} \sqrt{2\pi/(\lambda t)}$ by completing
the square + translation invariance of \texttt{volume}; (2) deduce
\texttt{partition\_pos} and \texttt{partition\_h\_zero}; (3) get
\texttt{partition\_hasDerivAt} from the derivative of the closed
form, and show the target integral is $0$ by oddness. \dots\ That is
both stronger and likely shorter than the dominated proof.}

The key insight is that for the harmonic Gibbs the partition
function is computable in closed form, so there is no need to invoke
the abstract differentiation theorem at all. The closed form is also
a stronger primitive --- it gives the value of $Z(h)$ for all $h$,
not just at $h = 0$ --- and is genuinely useful for downstream work
(e.g.\ moment-generating function calculations) in a way the bare
\texttt{GibbsRegularity} instance is not.

\paragraph{The arithmetic correction.} GPT also caught a minor
arithmetic error in the pre-deliberation Young bound: Claude wrote
$(\lambda/2)x^2 + h\,x \ge (\lambda/4)x^2 - h^2/(2\lambda)$ when the
correct bound is $\ge (\lambda/4)x^2 - h^2/\lambda$ (factor of two
off in the second term). This bound becomes moot under the
closed-form route, but is worth recording for any future tide that
does \emph{do} dominated convergence on a non-Gaussian potential
where the Young inequality doesn't dissolve into a closed form.

\paragraph{The \texttt{conv\_lhs => ext x; rw [\dots]} non-tactic.}
A first attempt at rewriting the integrand inside an integral via
\texttt{conv\_lhs} failed: \texttt{ext x} doesn't dig into integral
expressions, only into function-typed targets. The error was
``\texttt{invalid 'ext' conv tactic, function or arrow expected:
\(\int (x : \mathbb{R}), \dots : \mathbb{R}\)}''. The fix was to
prove the integrand-equality at the lambda level
($f = g$ as functions) and then \texttt{rw} the resulting integral
equality. Roughly:

\begin{verbatim}
have h_integrand_eq : (fun x => ...) = (fun x => ...) := by
  funext x; <ring/exp_add manipulation>
rw [show (integral x, ...) = (integral x, ...)
    from by rw [h_integrand_eq]]
\end{verbatim}

This is a recurring Lean idiom for ``rewrite the integrand of an
integral'' and is worth landing in the laplace
\texttt{CLAUDE.md}.\footnote{It is closely related to the
``integrand-lambda gymnastics'' issue recorded in the
anharmonic-$\kappa_3$ retrospective: when an integral's integrand is
a complicated expression, lambda-form rewriting plus
\texttt{integral\_congr} (or just substitution into the integral
binding) is more robust than \texttt{conv\_lhs => ext x}.}

\paragraph{\texttt{hasDerivAt\_pow} type ascription.} Mathlib's
\texttt{hasDerivAt\_pow} is generic over the field: it takes a
natural-number exponent and a base in any normed field $\mathbb{K}$.
Calling it as \texttt{hasDerivAt\_pow 2 0} fails because Lean tries
to infer $0 : \mathbb{N}$ first. The fix is to ascribe the base
type: write \texttt{hasDerivAt\_pow 2 (0 : R)}, where \texttt{R} is
the real-number type. Same pattern recurred several times in the
closed-form derivative proof; documented in \texttt{CLAUDE.md} as a
known type-ascription friction.

\paragraph{\texttt{HasDerivAt.exp} for $\exp \circ g$.} The natural
form of the chain rule for $\exp$ is \texttt{HasDerivAt.exp}, not
\texttt{Real.hasDerivAt\_exp.comp}. The latter doesn't typecheck
because \texttt{Real.hasDerivAt\_exp} is a $\forall x, \texttt{HasDerivAt}\dots$
statement (a family of facts at each point), not a single
\texttt{HasDerivAt} fact. The correct usage is
$\texttt{h\_g.exp} : \texttt{HasDerivAt}\,(\exp \circ g)\, (\exp(g(0)) \cdot g'(0))\,0$
when $\texttt{h\_g} : \texttt{HasDerivAt}\,g\,g'(0)\,0$. Two minutes
of search-and-replace once located.

\paragraph{Final \texttt{ring} closing a closed goal.} At the end of
\texttt{harmonic\_linear\_partition\_eq}, the proof structure
\texttt{congr 1; field\_simp; ring} produced ``No goals to be
solved''. The \texttt{field\_simp} had already closed the goal, so
the trailing \texttt{ring} was redundant. Removed it. (This is the
``\texttt{field\_simp} sometimes closes the goal, sometimes leaves
residue for \texttt{ring}'' inconsistency already documented in the
laplace \texttt{CLAUDE.md}.)

\section{What was Mathlib, what was new}

\paragraph{Mathlib provided.}
\begin{itemize}
\item \texttt{integral\_gaussian}: $\int_{\mathbb{R}} \exp(-b\,x^2)\,dx
= \sqrt{\pi / b}$ for any $b \in \mathbb{R}$.
\item Translation invariance of Lebesgue measure on
$\mathbb{R}$.\footnote{The additive analogue (via
\texttt{to\_additive}) of \texttt{integral\_mul\_right\_eq\_self} in
\texttt{MeasureTheory.Group.Integral}.}
\item \texttt{integral\_const\_mul} for pulling a constant out of an
integral.
\item \texttt{HasDerivAt.const\_mul}, \texttt{HasDerivAt.div\_const},
\texttt{HasDerivAt.exp}, \texttt{HasDerivAt.mul\_const},
\texttt{hasDerivAt\_pow}: standard \texttt{HasDerivAt}-algebra glue.
\item \texttt{Real.sqrt\_pos}, \texttt{Real.exp\_zero},
\texttt{Real.exp\_add}, \texttt{positivity}, \texttt{ring},
\texttt{field\_simp}, \texttt{funext}.
\end{itemize}

\paragraph{Threepoint provided.}
The \texttt{GibbsRegularity} structure itself.

\paragraph{What was new.}
\begin{itemize}
\item \texttt{harmonic\_square\_completion} (private): pointwise
identity. Three lines.
\item \texttt{exp\_neg\_t\_harmonic\_linear\_eq} (private): pointwise
factorisation of the Boltzmann factor. Six lines.
\item \texttt{harmonic\_linear\_partition\_eq} (public, headline
helper): the closed-form partition for all $h$. Twenty-five lines.
\item \texttt{harmonic\_partition\_eq} (public): $h = 0$
specialisation, $Z(0) = \sqrt{2\pi/(\lambda t)}$. Twenty lines.
\item \texttt{harmonic\_partition\_pos} (public): the first
\texttt{GibbsRegularity} field. Five lines.
\item \texttt{integral\_eq\_zero\_of\_odd} (private): mirror of the
helper from the harmonic-$\kappa_3$ tide. Eight lines. (Worth
promoting to \texttt{lean-common} on this second use.)
\item \texttt{harmonic\_integral\_neg\_t\_x\_exp\_neg\_eq\_zero}
(public): the right-hand-side derivative integral, zero by parity.
Six lines.
\item \texttt{closed\_form\_hasDerivAt\_zero} (private): derivative
of the closed-form partition at $h = 0$ is zero. Twenty-five lines
including the chain-rule plumbing.
\item the headline structure instance\footnote{Full name:
\texttt{Threepoint.harmonic\_id\_gibbsRegularity}.} (public): twenty
lines.
\end{itemize}

The mathematical content is the closed form (an elementary
calculation) and the parity-vanishing (already proven for the same
weighted integrand in Tide~5, but on a different exponent
parameterisation, so re-derived). Everything else is plumbing.

\section{Lessons}

\begin{itemize}
\item \textbf{Look for closed forms before the abstract machinery.}
The pre-deliberation plan was to apply Mathlib's
dominated-convergence theorem; GPT's pivot to the closed-form route
saved an estimated 80--150 lines of plumbing and produced a stronger
primitive.
This generalises beyond the harmonic case: \emph{whenever the
partition or the moment integral is expressible in closed form, the
closed form is the right primitive theorem to expose}, and the
abstract regularity is a corollary. The dominated-convergence
machinery is reserved for cases where no closed form is available.

\item \textbf{The \texttt{conv => ext x} idiom does not work for
integral integrands.} \texttt{conv\_lhs => ext x; rw [\dots]} fails
inside integral expressions because \texttt{ext} expects a
function-typed target. Workaround: prove the integrand-equality at
the function level (via \texttt{funext}-and-manipulation) and
\texttt{rw} the resulting integral-level equality. Worth landing in
\texttt{CLAUDE.md} as a recurring idiom alongside the
``Pi.add vs single-lambda'' note from earlier tides.

\item \textbf{\texttt{HasDerivAt.exp}, not
\texttt{Real.hasDerivAt\_exp.comp}.} The chain rule for $\exp \circ g$
is \texttt{h\_g.exp}, with the value derivative
$\exp(g(x_0)) \cdot g'(x_0)$. The pointwise
\texttt{Real.hasDerivAt\_exp} is a $\forall x$ statement and does not
compose directly. Document.

\item \textbf{The closed-form route's stronger primitive is generic.}
The route generalises to any potential of the form
$L(x) = (\lambda/2) x^2 + Q(x)$ where $Q$ admits a closed-form
$\int e^{-tQ(x)} g(x)\,dx$ in terms of standard integrals --- in
particular to any potential whose Gibbs measure remains Gaussian after
linear perturbation, or where the perturbation can be absorbed by
substitution. For non-quadratic potentials (the anharmonic case from
laplace, the bounded-prior cases, the sextic and pure-quartic from
recent tides), no closed form for $Z(h)$ exists and the
dominated-convergence route is the right one. Future tides that need
\texttt{GibbsRegularity} for those cases will pay the dominated
convergence cost.

\item \textbf{Numerical sanity is sometimes not feasible, and that's
fine.} This tide's target is structural (an instance of
\texttt{GibbsRegularity}, a \texttt{Prop}-valued structure), not a
numerical asymptotic. The skill's discipline says to record
``\textsection Numerical check: not feasible: <reason>'' rather than
fake one; that is the right call here. The closed-form derivation is
self-verifying via \texttt{integral\_gaussian}, but it isn't a
numerical sanity check in the same sense as the
anharmonic-$\kappa_3$ tide's \texttt{scipy.integrate.quad}
comparison.
\end{itemize}

\section{Follow-ups}

\begin{itemize}
\item \emph{Cross-susceptibility derivative theorems for the
harmonic case.} With the regularity instance now in hand, the FDT
and cross-susceptibility derivatives specialise concretely. The
cross-susceptibility derivative is particularly satisfying: by Tide
5, $\kappa_3(\varphi, x, B) = 0$ for the harmonic case, so
\[
  \partial_h \mathrm{Cov}_h(\varphi, B)\big|_{h=0}
    = -t \cdot \kappa_3(\varphi, x, B) = 0
\]
for any $\varphi, B$ for which the relevant
\texttt{GibbsObservable} hypotheses hold. A natural one-tide step
would package this as a stand-alone theorem in laplace.

\item \emph{\texttt{GibbsObservable} instances for monomials.} The
companion structure \texttt{GibbsObservable} (the $\varphi$-specific
analytic content) is also still parameterised. Concrete instances
for $\varphi(x) = x^k$ at the harmonic potential would follow the
same closed-form route: the perturbed numerator
$\int x^k e^{-t((\lambda/2)x^2 + h\,x)}\,dx$ has a closed form (it's
a derivative of $Z(h)$ in some other parameter, or a moment of the
shifted Gaussian).

\item \emph{Promotion to \texttt{lean-common}.}
\texttt{integral\_eq\_zero\_of\_odd} has now appeared in two
files\footnote{Once in \texttt{Threepoint/Harmonic.lean}, once
here.} and is a self-contained generic statement. Worth promoting to
\texttt{lean-common} as a small \texttt{Parity.lean} module on the
next tide that needs it.

\item \emph{Stronger primitive: the moment-generating function.} The
closed form $Z(h) = e^{t h^2/(2\lambda)} \cdot \sqrt{2\pi/(\lambda
t)}$ is, up to the $\sqrt{2\pi/(\lambda t)}$ normalisation, the
\emph{moment-generating function} of the Gaussian with variance
$1/(\lambda t)$. Future tides that need higher harmonic moments can
read them off from $Z(h)$ via differentiation; this is the natural
generic primitive to expose.

\item \emph{Generalise to arbitrary at-most-linear $A$ (Candidate
C).} The dominated-convergence route would handle this; the
closed-form route does not, since the partition is no longer
Gaussian for nonlinear $A$. A future tide could implement the
dominated-convergence route for the harmonic + general $A$ case as
a stepping-stone for non-Gaussian potentials.

\item \emph{Add the new gotchas to \texttt{laplace/CLAUDE.md}.}
Three items: the \texttt{conv\_lhs => ext x} non-applicability for
integrals, the \texttt{HasDerivAt.exp}-not-\texttt{.comp} pattern,
and the closed-form-first heuristic for \texttt{GibbsRegularity}
instances.
\end{itemize}

\section*{Acknowledgements}

GPT-5.5~Pro contributed the deliberation-stage advisory at Step~2,
including the strategic pivot from dominated convergence to the
closed-form route, the corrected Young bound, and the verified
suggestion to look for \texttt{integral\_add\_right\_eq\_self} and
\texttt{integral\_gaussian} as the load-bearing Mathlib calls. The
full consult and the tide log itself are preserved in the patterning
project's tide-log directory.\footnote{\texttt{sri/projects/patterning/tide-log/}\allowbreak\texttt{2026-05-06-tide-harmonic-gibbsregularity.md}
and \texttt{gpt55\_tide\_harmonic\_gibbsregularity\_v1.md}.}

The seabed for this tide spans threepoint's
\texttt{GibbsRegularity} structure (Chris Elliott's earlier scaffold
in \texttt{Threepoint/CrossSusceptibility.lean}), Mathlib's Gaussian
integral and translation-invariance machinery, and the cross-repo
dependency added by Tide~9 (laplace
\texttt{[[require]]}-ing threepoint at commit \texttt{40398d8}).
The present tide is the first to actually \emph{instantiate}
threepoint's regularity bundle on a concrete potential, completing
the loop that Tides 5 and 9 deliberately deferred.

\end{document}
